% Generated draft: requires mathematical and source review.
\documentclass[11pt]{article}
\usepackage[margin=25mm]{geometry}
\usepackage{fontspec}
\setmainfont{Latin Modern Roman}
\usepackage{amsmath,amssymb,mathtools,unicode-math}
\setmathfont{Latin Modern Math}
\usepackage{enumitem,needspace,xurl,hyperref}
\setlength{\parindent}{0pt}
\setlength{\parskip}{5pt}
\setlength{\emergencystretch}{3em}
\newcommand{\N}{\mathbb{N}}
\newcommand{\Z}{\mathbb{Z}}
\newcommand{\Q}{\mathbb{Q}}
\newcommand{\R}{\mathbb{R}}
\newcommand{\C}{\mathbb{C}}
\newcommand{\F}{\mathbb{F}}
\newcommand{\A}{\mathbb{A}}
\newcommand{\PP}{\mathbb{P}}
\newcommand{\E}{\mathbb{E}}
\newcommand{\eps}{\varepsilon}
\newcommand{\dd}{\,\mathrm{d}}
\begin{document}
\section*{UnsolvedMath ID: 701550}
\textbf{Source problem code:} catalogue entry
\subsection{Definitions and mathematical statement}
For a group \(G\) and a subgroup \(L\leq G\), define the core of \(L\) in \(G\) by
\[
L_G=\bigcap_{g\in G}gLg^{-1},
\]
the largest normal subgroup of \(G\) contained in \(L\). Write \(1\) for the trivial subgroup. Following the supplied transcription, for \(M\trianglelefteq G\), say that \(G\) \emph{nearly splits over \(M\)} if there exists \(N\leq G\) such that
\[
[G:N]=\infty,\qquad [G:MN]<\infty,\qquad (M\cap N)_G=1.
\]
The proposed question is whether, for every group \(G\), every prime \(p\), and every normal subgroup \(H\trianglelefteq G\) of order \(p\),
\[
\psi(G)\cap H=1
\quad\Longleftrightarrow\quad
G\text{ nearly splits over }H.
\]
\textbf{Unresolved qualifications.} The supplied record does not define \(\psi(G)\), so the left-hand side cannot be given a complete mathematical interpretation from the supplied evidence. Moreover, the two index requirements in the transcribed definition are incompatible when \(M\) is finite, hence in particular when \(|H|=p\). This is a formulation issue requiring source clarification; no correction to the definition or to the hypothesis on \(H\) is assumed here.
\subsection{Short English statement}
M. K. Azarian asks whether a normal subgroup \(H\) of prime order has trivial intersection with \(\psi(G)\) exactly when \(G\) nearly splits over \(H\), using the supplied definition above. The notation \(\psi(G)\) is undefined in the record, and the stated index conditions conflict with the finiteness of \(H\). The question therefore remains an unresolved transcription requiring review.
\subsection{Sources}
\begin{itemize}
\item Supplied local UnsolvedMath registry record, permanent ID \(701550\), exact source problem code \texttt{catalogue entry}, titled \emph{Kourovka Notebook Problem 19.2}; upstream revision \texttt{d371ee917ac9647fbb06cd6fc529deec6f3dda82}.
\item The record attributes the problem to M. K. Azarian, proposed in 2018, and cites \emph{The Kourovka Notebook (Unsolved Problems in Group Theory)}, 19th issue (2018), original problem number 19.2, together with the 21st issue edition (2026), identified as \texttt{arXiv:1401.0300v48}. The supplied collection description names E. I. Khukhro and V. D. Mazurov as editors.
\item The registry labels this problem \emph{open}. This is supplied source metadata, not an independently verified mathematical or literature assessment. The cited publication has not been consulted for this draft; no definition of \(\psi(G)\) or explanation of the index conflict is supplied.
\end{itemize}
\end{document}
